% from main.tex
%%%%%%%%%%  scan idx 417  |  folio 404  %%%%%%%%%%
\textbf{Corollaire 5.4.} La filtration induite par $K^{\bullet}(X)$ sur $K^{\bullet}(S)$ est identique à la filtration de $K^{\bullet}(S)$.

\textbf{Corollaire 5.5.} Soient $S$ un schéma quasi-compact, $E_1,\ldots,E_n$ des $\mathcal O_S$-Modules localement libres de type fini, $D(E_1)$ le schéma de drapeaux de $E_1$, $X=D(E_1)\times_S\cdots\times_S D(E_n)$. Alors l'homomorphisme canonique :
\[
  \Gr K^{\bullet}(S)\longrightarrow \Gr K^{\bullet}(X)
\]
est injectif.

Le raisonnement de 3.1 permet de se ramener à montrer la même propriété quand $X$ est un fibré projectif sur $S$. Ce cas résulte alors immédiatement de 5.4.

\textbf{Corollaire 5.6.} Avec les notations et les hypothèses de 5.1, $\Gr K^{\bullet}(X)$ est un $\Gr K^{\bullet}(S)$-module libre, admettant pour base $1,\bar x,\ldots,\bar x^r$. Pour tout $k$, on a :
\[
  \Gr^kK'(X)=\Gr^kK'(S)\oplus \Gr^{k-1}K^{\bullet}(S).\bar x\oplus\cdots\oplus \Gr^{k-r}K^{\bullet}(S).\bar x^r .
\]
Enfin, les éléments $1,\bar x,\ldots,\bar x^{r+1}$ sont liés par la relation :
\[
  \sum_{i=0}^{r+1}c^{r+1-i}(\check E)\bar x^i=0,
\]
où $c^k(\check E)$ est la $k$-ième classe de Chern de $\check E$ (V 6.7).

Conséquence immédiate de 1.1, 5.3, et (1.13.1).

\textbf{Corollaire 5.7.} Avec les notations et les hypothèses de 4.5, la $\Gr K^{\bullet}(S)$-algèbre $\Gr K^{\bullet}(X)$ est engendrée par les $\bar c_i^{(j)}$, où $\bar c_i^{(j)}$ est la classe dans $\Gr^iK'(X)$ de $\gamma^i(F_j-p_j)$. L'idéal des relations entre les $\bar c_i^{(j)}$ est engendré par les relations déduites de :
\[
  c(E-n)=\prod_{j=1}^{k}c(F_j-p_j),
\]
